%%
%% CSD Assignment 1 - Final Technical Report
%% ACM "acmart" manuscript style (compacted for <= 10 pages)
%%
\documentclass[manuscript,screen,review]{acmart}

\usepackage{booktabs}
\usepackage{tabularx}
\usepackage{amsmath}
\usepackage{amsfonts}
\usepackage{tikz}
\usetikzlibrary{shapes.geometric, arrows.meta, positioning, fit, backgrounds, calc}

% Compact float and list spacing for <= 10 pages target
\setlength{\textfloatsep}{8pt plus 2pt minus 2pt}
\setlength{\floatsep}{6pt plus 2pt minus 2pt}
\setlength{\intextsep}{6pt plus 2pt minus 2pt}

\AtBeginDocument{%
  \providecommand\BibTeX{{Bib\TeX}}}

\begin{document}

%% =================== TITLE & AUTHOR BLOCK =========================
\title{AI-based Village Pond Planning System: Automated Hydrological Modeling, Catchment Delineation, and Runoff Estimation for Rural Rainwater Harvesting}
\subtitle{CSD Assignment 1 -- Final Technical Report}

\author{Donga Shanmukha Siva Venkata Sai}
\affiliation{%
  \institution{Department of Computer Science and Engineering}
  \city{Hyderabad}
  \country{India}}
\email{shanmukhasivavenkatasaidonga@gmail.com}

\renewcommand{\shortauthors}{D. S. S. V. Sai}

%% =================== ABSTRACT ======================================
\begin{abstract}
Rural rainwater harvesting is vital for sustaining groundwater aquifers and rainfed agriculture in semi-arid regions. However, optimal farm-pond siting is impeded by expensive manual surveys, subjective field heuristics, and inaccessible desktop GIS software. In this work, we present the \emph{AI-based Village Pond Planning System}, an automated geospatial hydrology and decision-support framework integrated into a zero-install Web GIS application. Given contour vectors in KML or GeoJSON formats, the system reprojects coordinates to local Universal Transverse Mercator (UTM Zone 44N) metrics, interpolates a continuous Digital Elevation Model (DEM) via 2D Delaunay-based bivariate triangulation, and derives key terrain properties including slope and the Topographic Wetness Index (TWI). Hydrological flow is modeled using a vectorized D8 steepest-descent algorithm coupled with an inverted Compressed Sparse Row (CSR) graph traversal engine, delivering sub-millisecond ($<1\,\text{ms}$) recursive upstream watershed delineation for any user-selected pour point. Candidate pond sites are prioritized through Multi-Criteria Evaluation (MCE) combining flow accumulation, topographic depressions, slope thresholds, and spatial non-maximum suppression. Harvestable runoff volume is estimated via the empirical Rational Method ($V = C \times P \times A$). Built with a decoupled FastAPI ASGI backend and an interactive Leaflet frontend featuring real-time land-parcel drawing and dynamic hydrology visualization, the platform executes the end-to-end pipeline in $1.15\,\text{s}$ over a $55,000$-cell terrain grid. A real-world case study in Andhra Pradesh demonstrates the identification of an optimal $3,750\,\text{m}^3$ pond site served by a $3.245\,\text{ha}$ catchment yielding $10.79\,\text{ML}$ in annual runoff.
\end{abstract}

\keywords{Village Pond Planning, Geospatial Hydrology, Rainwater Harvesting, D8 Flow Routing, Catchment Delineation, Multi-Criteria Evaluation, Web GIS, REST APIs}

\maketitle

%% ==================== PROJECT ACCESS LINKS ===========================
\begin{center}
  \vspace{-3mm}
  \noindent\fbox{%
    \parbox{0.95\linewidth}{%
      \centering\small
      \textbf{GitHub Repository URL:} \url{https://github.com/DongaShanmukhaSivaVenkataSai/Pond-Catchment-Site-Selection-System} \\[2pt]
      \textbf{Final Working Front-end URL:} \url{http://10.1.75.51:5252}
    }%
  }
  \vspace{1mm}
\end{center}

%% =====================================================================
\section{Introduction}
\label{sec:intro}
Water scarcity severely constrains agricultural productivity across semi-arid regions of India, where over $60\%$ of cultivated land depends entirely on erratic monsoon rains. Farm ponds, percolation tanks, and check dams provide decentralized rainwater harvesting (RWH) infrastructure to recharge shallow aquifers and supply protective irrigation~\cite{critchley1991water}. However, many rural ponds fail or dry up prematurely due to poor spatial siting---frequently positioned in topographic saddles with deficient catchment areas or on steep slopes prone to severe siltation. 

This project develops an automated, web-accessible geospatial decision system that empowers village planners and agricultural officers to scientifically identify, evaluate, and locate rainwater harvesting ponds without requiring specialized GIS expertise.

\subsection{Motivation}
Manual theodolite and total-station field surveys are labor-intensive, slow, and cost-prohibitive for gram panchayats. Open-access satellite DEMs (e.g., SRTM 30m) are too coarse to resolve the meter-scale swales and field bunds that control agricultural runoff. While high-resolution elevation contours from UAV or cadastral surveys exist, field engineers lack automated tools to transform contour polylines into stream networks and catchment basins without complex desktop GIS packages (ArcGIS, QGIS). Our browser-based platform eliminates this bottleneck by unifying contour ingestion, DEM interpolation, D8 flow routing, and runoff estimation into a responsive, zero-install Web GIS application.

\subsection{Scope of the Project}
\textbf{In-Scope:} (1) Automated parsing of KML, KMZ, and GeoJSON contour vectors; (2) Conformal UTM metric coordinate reprojection; (3) 2D Delaunay-based DEM grid rasterization; (4) Terrain derivatives (slope, aspect, TWI); (5) Vectorized D8 flow routing and continuous flow accumulation; (6) MCE site selection with spatial Non-Maximum Suppression (NMS); (7) Sub-millisecond recursive upstream watershed delineation via inverted CSR graphs; (8) Rational Method runoff estimation ($V = C \cdot P \cdot A$); and (9) An interactive Leaflet Web GIS supporting user land-boundary polygon drawing and real-time metric updates.

\textbf{Out-of-Scope:} Sub-surface borehole lithological testing, structural civil design of masonry spillways or earthen bunds, dynamic hydrodynamic flood wave routing, and legal cadastral land-title dispute arbitration.

%% =====================================================================
\section{Problem Statement and Requirements}
\label{sec:requirements}
\textbf{Problem Statement:} \emph{Given elevation contour polylines and user-defined boundary constraints for a rural watershed, automatically identify optimal pond sites, delineate their contributing catchment areas, and compute harvestable runoff volume based on local precipitation and soil characteristics.}

Table~\ref{tab:requirements} maps each functional requirement to its implementation module.

\begin{table}[h]
  \small
  \caption{Mapping of functional requirements to implementation modules}
  \label{tab:requirements}
  \begin{tabular}{lp{8.0cm}}
    \toprule
    Requirement & Implemented in Module / File \\
    \midrule
    Satellite imagery display & \texttt{frontend/js/map.js} (CARTO Dark Matter \& Esri Satellite) \\
    Contour map visualization & \texttt{backend/app/services/kml\_parser.py}, \texttt{frontend/js/map.js} \\
    Available-land identification & \texttt{frontend/js/map.js} (Leaflet-Draw polygon \& Turf.js clipping) \\
    Catchment area estimation & \texttt{backend/app/services/catchment.py} (\texttt{UpstreamGraph} BFS) \\
    Historical rainfall query & \texttt{backend/app/services/hydrology.py}, \texttt{frontend/js/hydrology.js} \\
    Runoff volume estimation & \texttt{backend/app/services/hydrology.py}, \texttt{frontend/js/hydrology.js} \\
    Pond depth / storage recommendation & \texttt{backend/app/services/pond\_selector.py} (MCE sink scoring) \\
    Combined overlay / results view & \texttt{frontend/index.html}, \texttt{frontend/js/app.js}, \texttt{frontend/js/ui.js} \\
    \bottomrule
  \end{tabular}
\end{table}

\subsection{Non-Functional Requirements}
\begin{enumerate}
    \item \textbf{Performance:} Complete end-to-end pipeline execution on $50,000+$ contour points in $<2.0\,\text{s}$. On-the-fly catchment queries must execute in $<100\,\text{ms}$ for fluid map interaction.
    \item \textbf{Concurrency:} Dispatches CPU-bound numerical matrix operations to background thread pools via \texttt{run\_in\_executor}, preventing ASGI event loop starvation.
    \item \textbf{Usability:} Lightweight, zero-install Web GIS running in modern browsers with a total client bundle footprint $<1.5\,\text{MB}$ without heavy dependencies (e.g., React).
    \item \textbf{Resilience:} Stateless RESTful API with robust exception handling for boundary edge cases, concave geometries, and zero-slope sink depressions.
\end{enumerate}

%% =====================================================================
\section{System Architecture and High-Level Design}
\label{sec:architecture}
The platform is designed as an integrated, multi-tier geospatial system bridging an interactive browser-based Web GIS frontend with a high-performance Python ASGI backend. Figure~\ref{fig:architecture} presents an elaborated architectural tree diagram detailing the modular decomposition of the frontend and backend tiers, their constituent components, and the bidirectional communication bridge linking them.

\begin{figure*}[t]
  \centering
  \resizebox{0.98\linewidth}{!}{%
  \begin{tikzpicture}[
    font=\sffamily\footnotesize,
    node distance=4mm and 4mm,
    root/.style={draw=teal!90, fill=teal!15, rounded corners=4pt, line width=1pt, align=center, font=\bfseries\small, inner sep=5pt},
    tier/.style={draw=blue!80, fill=blue!10, rounded corners=3pt, line width=0.85pt, align=center, font=\bfseries\footnotesize, inner sep=4pt, text width=6.6cm},
    subtier/.style={draw=indigo!70, fill=indigo!6, rounded corners=3pt, line width=0.75pt, align=center, font=\bfseries\scriptsize, inner sep=3pt, text width=3.15cm},
    box/.style={draw=gray!70, fill=gray!4, rounded corners=2pt, align=left, font=\tiny, inner sep=3pt, text width=3.15cm},
    bridge/.style={draw=orange!90, fill=orange!12, rounded corners=4pt, line width=0.9pt, align=center, font=\bfseries\scriptsize, inner sep=4pt, text width=13.8cm},
    bridgebox/.style={draw=orange!75, fill=white, rounded corners=2pt, align=left, font=\tiny, inner sep=3pt, text width=13.8cm},
    conn/.style={-{Stealth[length=2mm]}, draw=black!75, line width=0.75pt},
    biconn/.style={{Stealth[length=2mm]}-{Stealth[length=2mm]}, draw=orange!90, line width=1pt}
  ]
    % Root
    \node[root] (root) {AI-based Village Pond Planning System: End-to-End Architecture};

    % Level 1: Frontend vs Backend
    \node[tier, below left=7mm and 4mm of root] (fe) {FRONTEND TIER (Web GIS Client)\\\normalfont\scriptsize Browser Runtime • Served at \texttt{http://10.1.75.51:5252}};
    \node[tier, below right=7mm and 4mm of root] (be) {BACKEND TIER (FastAPI ASGI Engine)\\\normalfont\scriptsize Python 3.10+ • Container Port \texttt{5000} (PID 832141)};

    \draw[conn] (root.south) -| (fe.north);
    \draw[conn] (root.south) -| (be.north);

    % Frontend Sub-tiers
    \node[subtier, below left=4mm and 1mm of fe.south] (fe_ui) {Map \& Visualization\\(\texttt{map.js}, \texttt{style.css})};
    \node[subtier, below right=4mm and 1mm of fe.south] (fe_logic) {Interaction \& Hydrology\\(\texttt{app.js}, \texttt{hydrology.js}, \texttt{ui.js})};

    \draw[conn] (fe.south) -| (fe_ui.north);
    \draw[conn] (fe.south) -| (fe_logic.north);

    \node[box, below=2mm of fe_ui] (fe_ui_box) {%
      $\bullet$ \textbf{Leaflet 1.9.4:} Map canvas engine\\
      $\bullet$ \textbf{CARTO Dark Matter:} API-keyed tiles\\
      $\bullet$ \textbf{Esri Satellite:} Ground-truth layer\\
      $\bullet$ \textbf{Leaflet.Draw:} Land parcel bounding\\
      $\bullet$ \textbf{Overlays:} Contours, streams, pond marker%
    };

    \node[box, below=2mm of fe_logic] (fe_logic_box) {%
      $\bullet$ \textbf{Turf.js:} Spatial boundary intersection\\
      $\bullet$ \textbf{Hydrology Engine:} $V = C \cdot P \cdot A$ math\\
      $\bullet$ \textbf{Dynamic HUD:} Live volume \& depth card\\
      $\bullet$ \textbf{Interactive Sliders:} Rain ($P$), Soil ($C$)\\
      $\bullet$ \textbf{State Store:} Client GeoJSON caching%
    };

    % Backend Sub-tiers
    \node[subtier, below left=4mm and 1mm of be.south] (be_gw) {API Gateway \& Server\\(\texttt{main.py}, \texttt{routes.py}, \texttt{schemas.py})};
    \node[subtier, below right=4mm and 1mm of be.south] (be_eng) {Hydrological Domain Engine\\(\texttt{backend/app/services/*})};

    \draw[conn] (be.south) -| (be_gw.north);
    \draw[conn] (be.south) -| (be_eng.north);

    \node[box, below=2mm of be_gw] (be_gw_box) {%
      $\bullet$ \textbf{Uvicorn ASGI:} Async event loop\\
      $\bullet$ \textbf{StaticFiles Mount:} Serves \texttt{/frontend} at \texttt{/}\\
      $\bullet$ \textbf{Pydantic Validation:} Request schemas\\
      $\bullet$ \textbf{ThreadPoolExecutor:} Offloads matrix math\\
      $\bullet$ \textbf{OpenAPI Docs:} \texttt{/docs} \& \texttt{/redoc} swagger%
    };

    \node[box, below=2mm of be_eng] (be_eng_box) {%
      $\bullet$ \textbf{kml\_parser / projection:} UTM 44N metric\\
      $\bullet$ \textbf{dem\_builder:} SciPy Delaunay grid interpolation\\
      $\bullet$ \textbf{terrain:} 2D Central-diff slope \& TWI\\
      $\bullet$ \textbf{hydrology:} Vectorized D8 flow routing\\
      $\bullet$ \textbf{catchment:} Inverted CSR \texttt{UpstreamGraph}\\
      $\bullet$ \textbf{pond\_selector:} MCE scoring \& spatial NMS%
    };

    % Connection Bridge below both
    \node[bridge, below=20mm of $(fe_ui_box.south)!0.5!(be_eng_box.south)$] (bridge_title) {FRONTEND $\longleftrightarrow$ BACKEND CONNECTION \& NETWORK BRIDGE};

    \node[bridgebox, below=1.5mm of bridge_title] (bridge_details) {%
      $\bullet$ \textbf{Port Forwarding Pipeline:} External Access \texttt{http://10.1.75.51:5252} $\longleftrightarrow$ Fixed Forwarding $\longleftrightarrow$ Container Port \texttt{5000} (Uvicorn ASGI server).\\
      $\bullet$ \textbf{Client API Connector (\texttt{frontend/js/api.js}):} Asynchronous \texttt{fetch()} wrapper auto-detecting \texttt{window.location.origin}; transparently routes API calls.\\
      $\bullet$ \textbf{Request Dataflow:} Multi-part file upload (\texttt{POST /api/v1/analyze} with \texttt{contours\_1m.kml}) and JSON coordinate payloads (\texttt{POST /api/v1/catchment/custom}).\\
      $\bullet$ \textbf{Response Contracts:} Standards-compliant GeoJSON FeatureCollections (stream polylines, candidate points, catchment polygons) + JSON summary metrics.%
    };

    % Connections to bridge
    \draw[biconn] (fe_logic_box.south) -- ++(0,-0.3) -| ([xshift=-3.5cm]bridge_title.north);
    \draw[biconn] (be_gw_box.south) -- ++(0,-0.3) -| ([xshift=3.5cm]bridge_title.north);

  \end{tikzpicture}%
  }
  \caption{Elaborated architectural tree diagram illustrating the structural breakdown of the Frontend Web GIS tier, Backend FastAPI service tier, Hydrological Domain Engine, and the bidirectional connection bridge linking them.}
  \label{fig:architecture}
  \Description{Elaborated hierarchical architecture diagram showing frontend, backend, and their network connection bridge.}
\end{figure*}

\subsection{Technology Stack}
\begin{itemize}
    \item \textbf{Backend:} \emph{FastAPI (Python 3.10+)}~\cite{fastapi2018} on the \emph{Uvicorn ASGI} server. Provides native asynchronous coroutines, high-throughput JSON serialization, and automatic OpenAPI schema generation.
    \item \textbf{Scientific \& Spatial Libraries:} \emph{NumPy} for vectorized array kernels; \emph{SciPy} (\texttt{griddata}) for 2D Delaunay DEM interpolation; \emph{PyProj} for EPSG:4326 to EPSG:32644 metric projection; and \emph{Shapely} for topological polygon unions and vector geometry serialization.
    \item \textbf{Frontend:} \emph{Vanilla JavaScript (ES6+), HTML5, CSS3}, and \emph{Leaflet 1.9.4}~\cite{leaflet2010} with \emph{Leaflet.Draw} and \emph{Turf.js}. Avoids heavy JavaScript framework bloat to deliver sub-second initial browser loading.
    \item \textbf{Basemaps:} Authenticated \emph{CARTO Dark Matter} for high-contrast thematic vector visibility, paired with \emph{Esri World Imagery} satellite tiles.
\end{itemize}

%% =====================================================================
\section{Methodology}
\label{sec:methodology}

\subsection{Terrain and Elevation Analysis}
Raw contour polylines $\mathcal{L} = \{L_i\}_{i=1}^k$ with elevation attributes $z_i$ are reprojected from geodetic coordinates $(\lambda, \phi)$ to conformal metric UTM Zone 44N coordinates:
\begin{equation}
    (x_j, y_j) = \mathcal{T}_{\text{EPSG:4326} \to \text{EPSG:32644}}(\lambda_j, \phi_j)
\end{equation}
A 3D point cloud $\mathcal{P} = \{(x_m, y_m, z_m)\}$ is rasterized onto a uniform 2D grid ($x_c \in [x_{\min}, x_{\max}]$, $y_r \in [y_{\min}, y_{\max}]$) with resolution $\delta = 1.0\,\text{m}$. Elevations are derived via 2D piecewise linear barycentric interpolation across the Delaunay triangulation of $\mathcal{P}$, with nearest-neighbor padding outside the convex hull:
\begin{equation}
    \text{DEM}(r, c) = \begin{cases}
        z_{\text{linear}}(x_c, y_r), & \text{if } (x_c, y_r) \in \text{Conv}(\mathcal{P}) \\
        z_{\text{nearest}}(x_c, y_r), & \text{otherwise}
    \end{cases}
\end{equation}
Terrain slope $\theta(r, c)$ in degrees is derived using central finite differences:
\begin{equation}
    p = \frac{\text{DEM}(r, c+1) - \text{DEM}(r, c-1)}{2\delta}, \quad
    q = \frac{\text{DEM}(r+1, c) - \text{DEM}(r-1, c)}{2\delta}
\end{equation}
\begin{equation}
    \theta(r, c) = \arctan\left(\sqrt{p^2 + q^2}\right) \times \frac{180}{\pi}
\end{equation}
The Topographic Wetness Index (TWI)~\cite{beven1979physically} characterizes moisture accumulation tendencies:
\begin{equation}
    \text{TWI}(r, c) = \ln\left( \frac{\alpha}{\tan(\theta \cdot \frac{\pi}{180}) + 10^{-5}} \right)
\end{equation}
where $\alpha$ is the specific upslope contributing area per unit contour width.

\subsection{Catchment Area Delineation}
Hydrological flow routing utilizes the discrete D8 single-flow-direction algorithm~\cite{tarboton1991new}, directing runoff from cell $(r, c)$ to the steepest downhill neighbor among the 8 Moore adjacent cells:
\begin{equation}
    k^* = \arg\max_{k \in \{0, \dots, 7\}} \frac{\text{DEM}(r, c) - \text{DEM}(r + \Delta r_k, c + \Delta c_k)}{d_k}
\end{equation}
where $d_k = \delta$ for orthogonal neighbors and $d_k = \delta \sqrt{2}$ for diagonal neighbors.

To achieve real-time catchment queries, we construct an inverted flow adjacency graph (\texttt{UpstreamGraph}) during initialization. Source-target flow links are packed into flattened Compressed Sparse Row (CSR) arrays indexed by cell IDs. Extracting the upstream catchment basin $\mathcal{W}(c_{\text{pour}})$ requires only an elementary Breadth-First Search (BFS) from the pour point:
\begin{equation}
    \mathcal{W}(c_{\text{pour}}) = \{c_{\text{pour}}\} \cup \bigcup_{v \in \text{upstream}(u)} \mathcal{W}(v)
\end{equation}
This graph traversal completes in $<1\,\text{ms}$, allowing users to click anywhere on the map and see immediate watershed boundary polygon updates.

\subsection{Pond Siting and Runoff Sizing}
Cells with sufficient drainage accumulation ($\text{Acc} \ge \tau_{\text{flow}}$) and gentle slopes ($\theta \le 10^\circ$) are evaluated using Multi-Criteria Evaluation (MCE):
\begin{equation}
    S(r, c) = 0.60 \cdot \widehat{\text{Acc}}(r, c) + 0.30 \cdot (1 - \widehat{\theta}(r, c)) + 0.10 \cdot (1 - \widehat{z}(r, c))
\end{equation}
where $\widehat{\text{Acc}}$, $\widehat{\theta}$, and $\widehat{z}$ are min-max normalized values. Spatial Non-Maximum Suppression (NMS) enforces a minimum separation distance $d_{\min} \ge 100\,\text{m}$ between top candidates.

Annual harvestable runoff is computed using the empirical Rational Method~\cite{chow1988applied}:
\begin{equation}
    V_{\text{runoff}} = C \times \left(\frac{P}{1000}\right) \times A_{\text{catchment}}
\end{equation}
where $V$ is volume ($\text{m}^3$), $C \in [0.15, 0.50]$ is the runoff coefficient, $P$ is annual rainfall ($\text{mm}$), and $A$ is the catchment area ($\text{m}^2$).

%% =====================================================================
\section{Implementation}
\label{sec:implementation}

\subsection{Backend and API Design}
The service exposes clean RESTful endpoints validated with Pydantic schemas:
(1) \texttt{POST /api/v1/analyze}: Full pipeline execution returning GeoJSON streams, pond candidates, and catchments;
(2) \texttt{POST /api/v1/catchment/custom}: On-the-fly sub-millisecond watershed extraction for user-clicked coordinates;
(3) \texttt{POST /api/v1/dem}: Direct DEM raster matrix generation; and
(4) \texttt{POST /api/v1/export/csv}: Tabular CSV export of candidate metrics.
CPU-bound matrix routines are offloaded via \texttt{asyncio.get\_event\_loop().run\_in\_executor(None, ...)}, preventing ASGI event loop blocking.

\subsection{Frontend and Visualization}
Built with Leaflet 1.9.4, the client provides:
(1) Dual CARTO Dark Matter and Esri Satellite basemap toggles;
(2) Interactive land selection using Leaflet.Draw to restrict candidates within user property lines;
(3) Dynamic layer controls for contours, streams, pond markers, and watershed polygons; and
(4) A floating HUD displaying catchment area, expected runoff volume, recommended pond depth, and inflow safety ratio. Dynamic sliders permit instant recalculation of rainfall and runoff coefficients in the browser.

\begin{figure}[h]
  \centering
  \fbox{
    \begin{minipage}{0.88\linewidth}
      \centering
      \vspace{2pt}
      \textbf{[ Web GIS Interface Overview ]} \\
      \scriptsize
      $\bullet$ Basemaps: CARTO Dark Matter / Esri Satellite \quad $\bullet$ Contours: 1m Polylines \\
      $\bullet$ Selection: User-Drawn Land Polygon \quad $\bullet$ Catchment: $3.25\,\text{ha}$ Watershed \\
      $\bullet$ Recommended Pond: 2.5m Depth, $3,750\,\text{m}^3$ \quad $\bullet$ Inflow Volume: $10.79\,\text{ML}$ \\
      \vspace{2pt}
    \end{minipage}
  }
  \caption{Application interface overview showing results overlay and real-time hydrological HUD.}
  \label{fig:ui}
  \Description{Overview of the Leaflet Web GIS showing map overlays and metrics.}
\end{figure}

\subsection{Data Management}
The architecture is completely stateless. Uploaded contours and intermediate matrices are processed in-memory and cached by dataset hash tokens, eliminating persistent database latency while allowing export to GeoJSON, CSV, and GeoPackage formats.

%% =====================================================================
\section{CSD Themes and Topics Applied in the Project}
\label{sec:csd-themes}
Table~\ref{tab:csd-themes} maps core Computer System Design (CSD) principles to concrete engineering implementations within the project.

\begin{table*}[t]
  \footnotesize
  \caption{Mapping of core CSD themes and topics to system implementation}
  \label{tab:csd-themes}
  \begin{tabular}{p{3.2cm}p{4.8cm}p{6.8cm}}
    \toprule
    CSD Theme/Topic & Where used in the project & Justification / design rationale \\
    \midrule
    API Design (REST) & \texttt{backend/app/routes.py} (\texttt{/api/v1/analyze}, \texttt{/catchment}) & Versioned REST endpoints with strict Pydantic schemas, HTTP status codes, and GeoJSON contracts. \\
    Reverse Proxy / Forwarding & Deployment architecture ($5252 \to 5000$) & Decouples internal container port from external access point; permits horizontal worker scaling. \\
    Caching & In-memory CSR flow graph memoization & Retains pre-computed \texttt{UpstreamGraph} in memory, reducing catchment queries from $700\,\text{ms}$ to $<1\,\text{ms}$. \\
    Database / Spatial Indexing & CSR adjacency indexing and Shapely STRtree & Replaces $O(M \cdot H)$ pixel iteration; enables sub-millisecond point-in-polygon clipping. \\
    Concurrency / Async & \texttt{asyncio} loop with \texttt{ThreadPoolExecutor} & Prevents CPU-bound NumPy/SciPy operations from blocking the ASGI loop, preserving client throughput. \\
    Microservices vs. Monolith & Modular service-oriented monolith & Eliminates distributed network overhead while maintaining clean module separation for parsing, DEM, and hydrology. \\
    Design Patterns & Pipeline and Strategy patterns & Pipeline pattern chains data transformations (KML $\to$ DEM $\to$ MCE); Strategy pattern enables swappable interpolators. \\
    Security & CARTO API token parameterization \& CORS & Secures commercial basemap access via token forwarding; enforces strict CORS origins against CSRF. \\
    Containerization & Linux container deployment with daemons & Encapsulates GDAL/PROJ dependencies in a reproducible environment deployed with background daemons. \\
    Error Handling \& Resilience & Geometric fallbacks \& Pydantic handlers & Gracefully handles concave contour boundaries, zero-division on flat terrain, and edge cases. \\
    Algorithms \& Complexity & Delaunay ($O(N\log N)$) \& CSR BFS ($O(K)$) & Replaces brute-force spatial scanning with topologically optimal graph algorithms. \\
    Testing Strategy & Pytest automated test suite & Verifies API status codes, coordinate reprojection, and conservation of mass in flow accumulation. \\
    Version Control / CI-CD & Git repository with branch isolation & Tracks functional commits, isolates sensitive credentials in \texttt{.gitignore}, and provides reproducible builds. \\
    \bottomrule
  \end{tabular}
\end{table*}

Key engineering rationales include: (1) \emph{Async Worker Offloading}: Decouples HTTP socket handling from heavy NumPy matrix computation; (2) \emph{Inverted CSR Flow Graph}: Enables instant BFS upstream watershed extraction; and (3) \emph{Zero-Install Web GIS}: Eliminates bulky desktop GIS installations for village administrators.

%% =====================================================================
\section{Results and Evaluation}
\label{sec:results}
The system was evaluated using a 1-meter contour elevation survey dataset (\texttt{contours\_1m.kml}) covering a rural micro-watershed in Andhra Pradesh, India ($17.65^\circ\,\text{N}, 82.95^\circ\,\text{E}$).

\subsection{Hydrological Analysis and Sizing}
The surveyed terrain spans $4.52\,\text{hectares}$ ($45,200\,\text{m}^2$) with elevations between $102.4\,\text{m}$ and $128.7\,\text{m}$ ($\Delta z = 26.3\,\text{m}$). Interpolation at $\delta = 1.0\,\text{m}$ produced a grid of $250 \times 220$ cells ($55,000$ elevation nodes). The top MCE pond candidate was identified at elevation $104.2\,\text{m}$ with a gentle $2.1^\circ$ slope, draining an upstream catchment of $32,450\,\text{m}^2$ ($3.245\,\text{ha}$, $71.8\%$ of surveyed area).

Assuming local annual rainfall of $P = 950\,\text{mm}$ and clay-loam runoff coefficient $C = 0.35$, the expected annual harvestable runoff is:
\begin{equation}
    V_{\text{runoff}} = 0.35 \times \left(\frac{950}{1000}\right) \times 32,450 = 10,789.6\,\text{m}^3 \approx 10.79\,\text{ML}
\end{equation}
A farm pond sized at $50\,\text{m} \times 30\,\text{m} \times 2.5\,\text{m}$ yields an effective storage capacity of $3,750\,\text{m}^3$. The inflow safety ratio is:
\begin{equation}
    R = \frac{10,789.6\,\text{m}^3}{3,750.0\,\text{m}^3} = 2.87
\end{equation}
This satisfies empirical design guidelines ($2.0 \le R \le 4.0$), ensuring the pond reliably fills even in below-average monsoon years, with excess directed through spillways.

\subsection{Performance Benchmarks}
Latency was benchmarked across 50 iterations on an Intel Xeon 2.20GHz, 16GB RAM instance. Table~\ref{tab:benchmarks} summarizes execution times.

\begin{table}[h]
  \small
  \caption{Performance benchmarks on a 55,000-cell terrain grid}
  \label{tab:benchmarks}
  \begin{tabular}{lrr}
    \toprule
    Pipeline Stage & Mean Latency & Share \\
    \midrule
    KML Ingestion \& Reprojection & $124\,\text{ms}$ & $10.7\%$ \\
    SciPy 2D Delaunay DEM Interpolation & $678\,\text{ms}$ & $58.8\%$ \\
    Terrain Slope, Curvature \& TWI & $62\,\text{ms}$ & $5.4\%$ \\
    Vectorized D8 Flow Routing & $196\,\text{ms}$ & $17.0\%$ \\
    MCE Pond Ranking \& Spatial NMS & $38\,\text{ms}$ & $3.3\%$ \\
    GeoJSON Vector Polygonization & $54\,\text{ms}$ & $4.7\%$ \\
    \midrule
    \textbf{Total Pipeline Execution} & \textbf{1,152\,ms} & \textbf{100.0\%} \\
    \bottomrule
    \emph{Dynamic Catchment Query (Inverted BFS)} & \emph{$0.85\,\text{ms}$} & \emph{Instantaneous} \\
    \bottomrule
  \end{tabular}
\end{table}

The pipeline executes in $1.15\,\text{s}$, well within the $2.0\,\text{s}$ target, with subsequent catchment queries completing in $<1\,\text{ms}$.

%% =====================================================================
\section{Discussion and Limitations}
\label{sec:discussion}
The platform demonstrates that automated hydrological modeling can replace manual field heuristics and eliminate expensive desktop GIS requirements. Key limitations include: (1) \emph{Contour Spacing}: Sparse contours in flat valleys can create minor interpolation stair-stepping requiring numerical smoothing; (2) \emph{Uniform Rainfall Assumption}: The Rational Method assumes steady rainfall intensity over the catchment; future iterations will incorporate the USDA SCS Curve Number method for storm event modeling; and (3) \emph{Geotechnical Verification}: The tool acts as a spatial suitability filter and must be paired with physical soil permeability testing to prevent excessive percolation losses.

%% =====================================================================
\section{AI Tool Usage Declaration}
\label{sec:ai-usage}
In compliance with the course Academic Integrity and LLM Usage Policy, Google Gemini / Antigravity was utilized for: (1) Scaffolding boilerplate FastAPI route handlers and Pydantic schemas; (2) Reference syntax for Leaflet.Draw and Turf.js event integration; and (3) LaTeX document typesetting. All core mathematical models, graph traversal algorithms, spatial calculations, and conclusions were verified, tested, and validated by the author.

%% =====================================================================
\bibliographystyle{ACM-Reference-Format}
\bibliography{sample-base}

%% =====================================================================
\appendix

\section{Source Code and Working Deployment}
\begin{itemize}
    \item \textbf{GitHub Repository URL:} \url{https://github.com/DongaShanmukhaSivaVenkataSai/Pond-Catchment-Site-Selection-System}
    \item \textbf{Final Working Front-end URL:} \url{http://10.1.75.51:5252}
\end{itemize}

\begin{verbatim}
Pond-Catchment-Site-Selection-System/
|-- backend/
|   |-- app/services/ (catchment, dem_builder, hydrology, pond_selector)
|   |-- app/routes.py & models/schemas.py
|   `-- main.py (FastAPI app & static frontend mount)
|-- frontend/ (index.html, css/style.css, js/app.js, map.js, api.js)
`-- requirements.txt & report/ (CSD_Assignment1_Report.tex)
\end{verbatim}

\textbf{Execution:} Clone repository, run \texttt{pip install -r requirements.txt}, and start server via \texttt{uvicorn backend.main:app --host 0.0.0.0 --port 5000 --reload}. Access UI at \texttt{http://localhost:5000}.

\end{document}
\endinput
